\documentclass[a4paper,11pt]{ltjsarticle}
\usepackage{base}
\title{}
\author{}
\date{}
\begin{document}
\begin{toi}
    \noindent 硫酸の工業的製法である\fbox{　A　}では，\underline{硫黄を酸化して\fbox{　B　}とし}$_{①}$，\underline{これを酸化して\fbox{　C　}する}$_{②}$．次いで，\underline{\fbox{　C　}を濃硫酸二吸収させ，その中の水と反応させて硫酸を得る．}$_{③}$
    \begin{itemize}
    \item [(1)]\fbox{　A　}から\fbox{　C　}に当てはまる語句や化学式を答えよ．
    \item [(2)]各反応①～③の化学反応式を書け．ただし，触媒が必要な反応ではそれも記せ．
    \item [(3)]この製法を用いて硫黄$9.6$kgを用いて硫酸を作る場合，硫黄をすべて硫酸に変えたと仮定すると，98$\%$の濃硫酸は何kg得られるか．
\end{itemize}
\end{toi}

\begin{toi}
硝酸は火薬や占領，医薬品の製造などに幅広く利用されている．工業的には，硝酸はアンモニアから製造される．製造工程は，はじめに，アンモニアと空気を白金触媒の存在科で反応させ，化合物Aと水を得る(反応①)．次いで，化合物Aを酸化して化合物Bを得る(反応②)．最後に，化合物Bと水を反応させて硝酸と化合物Aが得られる(反応③)．反応③で得られた化合物Aは，製造工程の一部に循環されるのがポイントである．
\begin{itemize}
    \item [(1)]上記で示した硝酸の工業的製法の名称を答えよ．
    \item [(2)]反応①から反応③を化学反応式で書け．
    \item [(3)]化合物Bは，反応③において次の(あ)〜(え)のどの働きをするか．最も適切なものを1つ選べ．\\
    (あ)酸化剤~(い)還元剤~(う)酸化剤でも還元剤でもある~(え)酸化剤でも還元剤でもない
    \item[(4)]反応①から反応③を1つにまとめた化学反応式を書け．
    \item[(5)]上記の方法で濃硝酸(質量パーセント濃度$69\%$，密度 $1.4$g/cm$^3$)1.0Lを製造するために必要なアンモニア(気体) の体積を標準状態で求めよ．ただし，$\text{H}=1.0,\text{N}=14,\text{O}=16$．
    \item[(6)]濃硝酸は銅と反応し，下方置換法により捕集される有色の気体を発生する．この反応の化学反応式を書け． 
\end{itemize}
\end{toi}
\begin{toi}
\begin{itemize}
    \item [(1)]アンモニアの工業的製法の名称を答え，それを化学反応式で表せ．ただし，触媒が必要な場合はそれも記せ．
        \item [(2)]炭酸ナトリウムの工業的製法の名称を答えよ．
               \item [(3)]炭酸ナトリウムの工業的製法において，炭酸水素ナトリウムを生成する反応を化学反応式で表せ．
                  \item [(4)]炭酸ナトリウムの工業的製法において，炭酸水素ナトリウムから炭酸ナトリウムを得る反応を化学反応式で表せ．
    \end{itemize}
\end{toi}
\end{document}